\documentclass[11pt]{article}
\usepackage[margin=1in]{geometry}
\usepackage{amsmath,amssymb,amsthm,mathtools,bm}
\usepackage[T1]{fontenc}
\usepackage{lmodern}
\usepackage[colorlinks=true,linkcolor=blue,citecolor=blue,urlcolor=blue]{hyperref}
\usepackage{microtype}
\newtheorem{theorem}{Theorem}[section]
\newtheorem{proposition}[theorem]{Proposition}
\newtheorem{lemma}[theorem]{Lemma}
\newtheorem{corollary}[theorem]{Corollary}
\theoremstyle{definition}
\newtheorem{definition}[theorem]{Definition}
\newtheorem{assumption}[theorem]{Assumption}
\theoremstyle{remark}
\newtheorem{remark}[theorem]{Remark}
\newcommand{\ii}{\mathrm i}
\newcommand{\dd}{\,\mathrm d}
\newcommand{\Tr}{\operatorname{Tr}}
\newcommand{\Ran}{\operatorname{Ran}}
\newcommand{\supp}{\operatorname{supp}}
\newcommand{\dist}{\operatorname{dist}}
\newcommand{\Spec}{\operatorname{Spec}}
\newcommand{\CAR}{\operatorname{CAR}}
\newcommand{\dG}{\mathrm d\Gamma}
\newcommand{\one}{\mathbf 1}
\newcommand{\NN}{\mathcal N}
\newcommand{\HH}{\mathcal H}
\newcommand{\FF}{\mathcal F}
\newcommand{\TT}{\mathbb T}
\newcommand{\up}{\uparrow}
\newcommand{\dn}{\downarrow}
\newcommand{\epsL}{O(L^{-\infty})}
\newcommand{\eqL}{\simeq}
\newcommand{\norm}[1]{\left\lVert#1\right\rVert}
\newcommand{\abs}[1]{\left|#1\right|}
\newcommand{\ip}[2]{\left\langle#1,#2\right\rangle}
\newcommand{\ad}{\operatorname{ad}}
\newcommand{\diam}{\operatorname{diam}}
\newcommand{\ind}{\operatorname{Ind}}
\DeclareMathOperator{\Var}{Var}
\title{A continuum charge-transport index for the lowest Landau level}
\author{}
\date{Working note, 6 November 2026}
\begin{document}
\maketitle
\begin{abstract}
We give a continuum version of the charge-transport argument of
Bachmann--Bols--De Roeck--Fraas, arXiv:1810.07351.
The arrangement of the argument is retained: spatial and charge data,
an index theorem, verification from a spectral gap, examples, and the
boundary-factorization proof.  A finite-rank formulation is proved as
well, using a determinant instead of the phase of a single state.
The lowest Landau level is included in the ordinary continuum fermion
space, and particles outside it are assigned a fixed positive energy.
Gaussian projection kernels then supply locality, while physical
regional particle numbers retain integer spectra.  The construction
uses no localized orthonormal basis and no discrete tight frame.
The low-energy sector is a band of \(d\) states with width at most
\(w\), separated from higher states by a uniform gap \(\gamma>2w\).
Exact degeneracy is not assumed; exponentially small splitting is
included.  For uniformly bounded local LLL field polynomials, a
particle-addition estimate preserves this band under the enlargement
when the relevant filling stays below saturation.  Neither positivity
nor zero ground energy is needed.  Positive parents whose lowest
energy is zero also admit the enlargement at saturation.
Magnetic translation gives the constraint
\(dn_\sigma/N_\phi\in\mathbb Z+O(L^{-\infty})\), where \(d\)
is the rank of this band in the selected spin-number sector.
At balanced rational filling \(\nu=p/q\) in lowest terms,
\(d\) is a multiple of \(2q/\gcd(p,2)\): this is \(2q\)
for odd \(p\), and \(q\) for even \(p\).
The same divisibility holds for the full low-energy band at fixed total
particle number when the Hamiltonian has spin-exchange symmetry.
Appendices derive the required kernel, propagation, and spectral-filter
estimates.
No gap at nonzero boundary twist or local accessibility of the first
excited state is assumed.
\end{abstract}
\tableofcontents

\section{Introduction}

The index in \cite{BBDF19} measures charge transported through one of
two distant boundaries by a charge-conserving unitary.  Its proof
uses a dressed regional charge, separation of the two boundaries,
and the periodicity of an integer-valued charge.  Those are the
ingredients we preserve.  For a low-energy sector of rank \(d>1\), the
appropriate quantity is an unnormalized trace.  We prove this
extension explicitly; the determinant formulation is motivated by
\cite{BBDF21}, but no theorem from either reference is used without
proof here.  As in the finite-rank formulation of \cite{BBDF21},
this sector may contain several distinct eigenvalues.  The gap used
in the proof is above the entire sector, rather than between its
individual levels.

There are two separate issues in transferring this argument to a
Landau level.  First, projecting a physical region onto the LLL does
not give a projection.  Its second quantization therefore need not
be an integer-valued charge.  Second, a spatially localized phase
rotation does not preserve the LLL.  We retain the full spatial
Hilbert space to resolve both issues, and penalize its additional
modes to preserve the desired low-energy spectrum.

The penalty is particularly simple for a positive parent Hamiltonian
with zero-energy ground states.  We also treat general self-adjoint,
number-conserving polynomials of uniformly bounded degree and
coefficients in localized LLL fields.  Averaging the energy of adding
a particle gives a uniform bound on the change of ground energy
between number sectors.  At fillings bounded away from saturation,
this proves that a sufficiently large, size-independent penalty
preserves the desired spectral band and its gap to higher states.
No gaps in the
smaller particle-number sectors are needed for that comparison.

Only magnetic translations at the original boundary conditions are
needed.  The auxiliary parameter in the index proof parametrizes
exponentials of dressed charges; it is not a family of Hamiltonians
whose spectral gap has to be established.

We write \(d\) for the low-energy band rank, reserving \(p,q\)
for the coprime numerator and denominator of the total filling.
The general spin-resolved constraint is derived first.  We then
specialize to balanced populations, discuss the numerator and
denominator parity, and prove a fixed-total-number version using
spin-exchange symmetry.  The general polynomial class covers balanced
fillings \(0<\nu<2\); the positive-parent case also covers the
saturated endpoint.  The rank \(d\) is fixed and the gap above the
band is uniform in the selected number sector.  The exact multiplicity
of the lowest eigenvalue may be smaller than \(d\).

\section{The index theorem and its LLL setting}

\subsection{Spatial setup, fermions, and charges}

Let
\[
 \TT_L=\mathbb R^2/(L_x\mathbb Z\oplus L_y\mathbb Z),\qquad
 L=\min(L_x,L_y),\qquad
 c_*\le L_x/L_y\le C_*,
\]
where \(L\to\infty\).  Work in units with magnetic length one, with
\[
 M=N_\phi=\frac{L_xL_y}{2\pi}\in\mathbb N.
\]
The one-particle space \(\mathfrak h\) is the space of square-integrable
sections of the magnetic line bundle.  In the gauge \(A=(0,x)\) its
sections obey
\begin{equation}\label{eq:bc}
 f(x+L_x,y)=e^{\ii L_xy}f(x,y),\qquad
 f(x,y+L_y)=f(x,y).
\end{equation}
Let \(\FF=\FF_{\mathrm a}(\mathfrak h\otimes\mathbb C^2)\).
The annihilation operators are antilinear in their argument and satisfy
\[
 \{a_\sigma(f),a_{\sigma'}^\dagger(g)\}
 =\delta_{\sigma\sigma'}\ip f g,\qquad
 \norm{a_\sigma(f)}=\norm f_2.
\]
Set \(N_\sigma=\dG_\sigma(1)\), and fix the sector
\[
 \HH_{\bm n}=\{\psi:N_\sigma\psi=n_\sigma\psi,\ \sigma=\up,\dn\},
 \qquad n_\up+n_\dn\le C L^2.
\]

A bounded operator is \emph{local in \(X\)} if it acts on
\(\FF_{\mathrm a}(L^2(X)\otimes\mathbb C^2)\), tensored with the
identity on the complement, with the usual fermionic grading.
We use only even operators, so operators in disjoint regions
commute.  An operator is \emph{neutral} if it commutes with both
particle numbers.  Locality is always defined in the ambient Fock
space, before restricting to \(\HH_{\bm n}\).

For a measurable region \(X\), define the physical spin-up charge
\[
 Q_X=\dG_\up(1_X).
\]
Although \(Q_X\) is unbounded on \(\FF\), its exponential is a
bounded local unitary, and
\begin{equation}\label{eq:integercharge}
 e^{2\pi\ii Q_X}=1,\qquad
 0\le Q_X|_{\HH_{\bm n}}\le n_\up.
\end{equation}
All charges commute.  Conjugation by \(e^{\ii sQ_X}\) multiplies the
spin-up test function by a phase of modulus one and preserves its
spatial support.

An estimate \(A\eqL B\) means
\(\norm{(A-B)|_{\HH_{\bm n}}}\le C_kL^{-k}\) for every \(k\).
When a bound is in ambient Fock-space norm we say so explicitly.
All constants may depend on the fixed density bound, aspect-ratio
bounds, interaction, and gap, but not on \(L\).
Polynomial factors in \(L\), and square roots, preserve an
\(\epsL\) error.

Take
\[
 \Gamma=[0,L_x/2)\times S_y^1,\qquad Q=Q_\Gamma.
\]
Its boundary components are denoted \(\partial_-\) and \(\partial_+\),
at \(x=0\) and \(x=L_x/2\), respectively.  Let \(S_\pm\) be the
strips of transverse radius \(L_x/8\) about these boundaries, and
let \(B_\pm\) be the narrower strips of radius \(L_x/16\).
Define
\[
 Q_\pm=Q_{\Gamma\cap S_\pm},\qquad Q_m=Q-Q_--Q_+.
\]
The two \(S_\pm\) are separated by a distance proportional to \(L\).

\subsection{Transport data and the theorem}

Let \(P=P_L\) be an orthogonal projection of fixed finite rank \(d\)
in \(\HH_{\bm n}\).  The following hypotheses isolate the algebraic
content of the index proof.

\begin{assumption}[Transport and subspace data]\label{ass:index}
\leavevmode
\begin{enumerate}
\item \(U\) is an ambient unitary preserving both particle numbers,
and \([U,P]\eqL0\).
\item There are neutral boundary transports \(T_\pm\), localized in
\(B_\pm\), with polynomial sector norms, such that
\[
 U^\dagger Q U-Q\eqL T_-+T_+.
\]
The operators \(Q_\pm+T_\pm\) have self-adjoint local realizations
and integer spectra:
\begin{equation}\label{eq:transportnormalization}
 e^{2\pi\ii(Q_\pm+T_\pm)}=1.
\end{equation}
Here and below unbounded local charges are used through their
bounded sector restrictions or their ambient unitary exponentials.
\item There are bounded, self-adjoint, neutral \(K_\pm\), localized
in \(B_\pm\), with polynomial ambient norms, such that
\[
 [Q-K_--K_+,P]\eqL0.
\]
There are also self-adjoint, neutral \(K_\pm^U\), localized in
\(B_\pm\), with polynomial ambient norms, satisfying
\[
 K_\pm^U\eqL U^\dagger K_\pm U.
\]
\item For bounded, neutral ambient operators \(A,B\) supported in
\(S_-,S_+\), respectively, there is a uniform estimate
\begin{equation}\label{eq:groundcluster}
 \norm{PA(1-P)BP}\le \norm A\,\norm B\,O(L^{-\infty}).
\end{equation}
\end{enumerate}
\end{assumption}

The normalization \eqref{eq:transportnormalization} is part of the
transport data.  We do not infer it from a tensor-factor argument
inside a fixed-particle-number sector.  It will hold exactly for
magnetic translation.  Spatial support in this assumption means
ambient support, not support of a compression to the LLL.

\begin{theorem}[Continuum transport index]\label{thm:index}
Under Assumption~\ref{ass:index},
\begin{equation}\label{eq:index}
 \dist\bigl(\Tr(PT_-),\mathbb Z\bigr)=O(L^{-\infty}).
\end{equation}
Consequently the normalized transport
\(\ind_P(U)=d^{-1}\Tr(PT_-)\) lies within \(O(L^{-\infty})\) of
\(d^{-1}\mathbb Z\).
\end{theorem}

The proof is in Section~\ref{sec:proof}.  For \(d=1\),
\eqref{eq:groundcluster} is ordinary connected clustering, and the
dressed charge has the vector spanning \(\Ran P\) as an approximate
eigenvector.
The proof then reduces to a scalar phase calculation, as in the
organization of \cite{BBDF19}.  Keeping \(d\) arbitrary requires only
replacing that phase by the determinant of a compressed unitary.
In particular, we never apply the rank-one theorem separately to
vectors in a higher-rank subspace.  No equality of energies within
\(\Ran P\) is required by this abstract theorem or its proof.

\subsection{The LLL and an exact enlargement}

Put
\[
 D_x=\partial_x,\quad D_y=\partial_y-\ii x,\quad
 D=\tfrac12(D_x-\ii D_y).
\]
Let \(\Pi\) be the orthogonal projection onto
\(\mathfrak h_0=\ker(D_x+\ii D_y)\), of dimension \(M\).
For \(R\in\TT_L\), its evaluation and derivative-evaluation vectors
are determined by
\[
 \ip{e_R}f=f(R),\qquad \ip{g_R}f=(Df)(R),\qquad f\in\mathfrak h_0.
\]
They belong to \(\mathfrak h_0\), even though \(Df\) need not.
Appendix~\ref{app:kernel} proves, uniformly in the torus,
\begin{equation}\label{eq:tails}
 \norm{e_R}+\norm{g_R}\le C,\qquad
 \norm{1_{\{d(\cdot,R)>r\}}e_R}
 +\norm{1_{\{d(\cdot,R)>r\}}g_R}\le Ce^{-cr^2},
 \qquad
 \int_{\TT_L}|e_R\rangle\langle e_R|\dd R=\Pi.
\end{equation}

\begin{definition}[Local LLL field polynomial]\label{def:local}
Let
\[
 H_{\mathrm{LLL}}=\int_{\TT_L}h_R\dd R
\]
where \(h_R=h_R^\dagger\) is a neutral polynomial in annihilators
and creators with test vectors \(e_R,g_R\).  The number of monomials,
their degrees, and their coefficients are uniformly bounded in \(R\)
and \(L\).  More generally, one may use a fixed finite family
\(f_{\alpha,R}\in\mathfrak h_0\) with uniform norm and Gaussian
\(L^2\)-tail bounds of the form \eqref{eq:tails}.  All these fields
are smeared by genuine one-particle vectors.

By averaging the polynomial over the two gauge rotations, one may
and will choose every monomial to have equal numbers of creators
and annihilators of each spin.  Thus the monomials are even and
neutral individually.  No sign condition on \(h_R\) is imposed.
\end{definition}

\begin{definition}[Positive local LLL parent]\label{def:parent}
A positive parent is the special case of Definition~\ref{def:local}
with
\[
 h_R=\sum_{\alpha=1}^{s}B_\alpha(R)^\dagger B_\alpha(R),
\]
where the \(B_\alpha(R)\) are polynomials in annihilators with
the same uniform bounds.  In applications of the positive-parent
gap comparison, the target sector is assumed to contain zero modes.
\end{definition}

\begin{assumption}[Isolated low-energy band]\label{ass:band}
In the selected all-LLL number sector, list the eigenvalues of
\(H_{\mathrm{LLL}}\) as \(E_1(L)\le E_2(L)\le\cdots\), counted
with multiplicity.  There are size-independent constants
\(d\in\mathbb N\), \(w\ge0\), and \(\gamma>2w\) such that
\begin{equation}\label{eq:bandhypothesis}
 w_L:=E_d(L)-E_1(L)\le w,\qquad
 \gamma_L:=E_{d+1}(L)-E_d(L)\ge\gamma.
\end{equation}
Let \(P\) be the rank-\(d\) spectral projection onto these lowest
\(d\) states.  The positive gap above \(E_d\) ensures that this
projection includes complete eigenspaces.  If the all-LLL sector
itself has dimension \(d\), use \(E_{d+1}=+\infty\).
\end{assumption}

We call \(\Ran P\) the low-energy band or sector.  In particular,
\(d\) is not assumed to be the multiplicity of \(E_1\).
The condition \(\gamma>2w\) is a convenient sufficient separation
for the filters below, matching the spectral-band formulation in
\cite{BBDF21}.  If \(w_L\to0\) and the gap above the band stays
positive, it holds for all sufficiently large sizes after choosing
a fixed \(w<\gamma/2\).  Exponential decay of \(w_L\) is more
than is needed.  An exactly degenerate ground space is the special
case \(w_L=0\).

For either interaction definition the same polynomial defines a bounded operator
on the full Fock space, acting as the LLL operator tensored with the
identity on spectator modes.  The CAR give \(\norm{h_R}\le C\)
and \(\norm{H_{\mathrm{LLL}}}\le CM\).  For fixed \(\mu>0\), set
\begin{equation}\label{eq:penalty}
 H_{\mathrm{ext}}=H_{\mathrm{LLL}}
       +\mu\sum_{\sigma=\up,\dn}\dG_\sigma(1-\Pi).
\end{equation}

\begin{lemma}[Positive-parent comparison]\label{lem:comparison}
For Definition~\ref{def:parent}, in a fixed spin-number sector
or a fixed-total-number sector, the spectrum of
\(H_{\mathrm{ext}}\) below \(\mu\), including multiplicities,
equals that of \(H_{\mathrm{LLL}}\) on the all-LLL sector below
\(\mu\).  In particular, suppose its lowest energy is zero and
its lowest \(d\) states have width at most \(w<\mu\) and gap
at least \(\gamma>0\) above them.  The enlargement preserves
this band, with the same eigenvalues and spectral projection,
and a gap above it of at least \(\min(\gamma,\mu-w)\).
For \(w=0\), this includes the exact-zero-mode comparison.
\end{lemma}
\begin{proof}
Decompose the one-particle space as
\(\mathfrak h_0\oplus\mathfrak h_0^\perp\) for each spin.
The interaction commutes with the number \(r\) of particles in
the second summand.  On that block the penalty is \(\mu r\) and
the interaction is nonnegative.  Blocks with \(r\ge1\) have
energies at least \(\mu\); the \(r=0\) block is exactly the
original LLL problem.  Positivity holds in every particle-number
block because of Definition~\ref{def:parent}.  The top of the
selected band is at most \(w\), giving the asserted gap above it.
\end{proof}

\subsection{Particle addition and removal of positivity}
\label{sec:addition}

We now prove a uniform comparison of ground energies for the general
class in Definition~\ref{def:local}.  Throughout the next two lemmas
we work on the finite-dimensional \emph{LLL Fock space}.  Write
\(H=H_{\mathrm{LLL}}\), and let \(E_0(\bm m)\) be its lowest
energy at spin populations \(\bm m=(m_\up,m_\dn)\), where
\(0\le m_\sigma\le M\).  Let \(E_{\mathrm{tot}}(m)\) be its
lowest energy at total population \(m\), allowing all spin sectors.

\begin{lemma}[Averaged particle addition]\label{lem:addition}
There is a constant \(C_0\), independent of \(M\) and the populations,
such that, whenever \(m_\sigma<M\),
\begin{equation}\label{eq:addspin}
 E_0(\bm m+\bm e_\sigma)-E_0(\bm m)
 \le\frac{C_0M}{M-m_\sigma}.
\end{equation}
Here \(\bm e_\sigma\) is the unit vector for the indicated spin.
For total populations \(0\le m<2M\), the same constant can be
chosen so that
\begin{equation}\label{eq:addtotal}
 E_{\mathrm{tot}}(m+1)-E_{\mathrm{tot}}(m)
 \le\frac{C_0M}{2M-m}.
\end{equation}
No spectral-gap assumption is made in this lemma.
\end{lemma}
\begin{proof}
Choose any orthonormal basis \(u_1,\ldots,u_M\) of
\(\mathfrak h_0\), and put \(a_{\sigma j}=a_\sigma(u_j)\).
Define
\[
 \mathcal C_\sigma(H)
   =\sum_{j=1}^M a_{\sigma j}[H,a_{\sigma j}^\dagger].
\]
We first prove the volume bound
\begin{equation}\label{eq:addcommutator}
 \norm{\mathcal C_\sigma(H)}\le C_0M.
\end{equation}
For a monomial \(F_1\cdots F_s\) of even degree in the local
fields, commuting a creator through the product gives exactly
\[
 [F_1\cdots F_s,a_{\sigma j}^\dagger]
 =\sum_{k=1}^s(-1)^{s-k}
 F_1\cdots F_{k-1}\{F_k,a_{\sigma j}^\dagger\}
 F_{k+1}\cdots F_s.
\]
The anticommutator vanishes unless \(F_k=a_\sigma(f)\), in
which case it is \(\ip f{u_j}\).  Multiplying on the left by
\(a_{\sigma j}\) and summing over \(j\) uses completeness only:
\[
 \sum_j\ip f{u_j}a_{\sigma j}=a_\sigma(f).
\]
Thus \(\mathcal C_\sigma(h_R)\) is a sum of a uniformly bounded
number of products of the same smeared fields.  Their norms are
uniformly bounded by the CAR and Definition~\ref{def:local}.
Integration over area \(2\pi M\) proves \eqref{eq:addcommutator}.
An orthonormal basis is used here only to take a complete sum;
no locality of these orbitals is asserted or needed.

On the LLL Fock space, \(\sum_j a_{\sigma j}a_{\sigma j}^\dagger
=M-N_\sigma\).  Hence
\begin{equation}\label{eq:addidentity}
 \sum_j a_{\sigma j}H a_{\sigma j}^\dagger
       =(M-N_\sigma)H+\mathcal C_\sigma(H).
\end{equation}
In particular, \(\mathcal C_\sigma(H)\) is self-adjoint, since
\([H,N_\sigma]=0\).
For a normalized ground vector \(\psi\) in sector \(\bm m\),
the trial vectors \(a_{\sigma j}^\dagger\psi\) lie in sector
\(\bm m+\bm e_\sigma\), and their squared norms sum to
\(M-m_\sigma\).  The variational principle and
\eqref{eq:addidentity} give
\begin{align*}
 (M-m_\sigma)E_0(\bm m+\bm e_\sigma)
 &\le\sum_j\ip{a_{\sigma j}^\dagger\psi}
                       {H a_{\sigma j}^\dagger\psi}\\
 &=(M-m_\sigma)E_0(\bm m)
       +\ip\psi{\mathcal C_\sigma(H)\psi}.
\end{align*}
Division by \(M-m_\sigma>0\) proves \eqref{eq:addspin}.
Individual trial vectors need not be nonzero or orthogonal.

For \eqref{eq:addtotal}, start instead from a total-\(m\)
ground vector and sum over both \(\sigma\) and \(j\).
The total trial norm is \(2M-m\), and the commutator correction
is \(\mathcal C_\up(H)+\mathcal C_\dn(H)\).  Increasing
\(C_0\) once to bound its norm by \(C_0M\) proves the claim.
\end{proof}

\begin{lemma}[Uniform energy bound for particle removal]
\label{lem:removal}
Fix \(\epsilon>0\), independent of size, and put
\(C_\epsilon=C_0/\epsilon\).
If \(n_\sigma\le(1-\epsilon)M\) for both spins, then
\begin{equation}\label{eq:removalspin}
 E_0(\bm n-\bm r)\ge E_0(\bm n)-C_\epsilon|\bm r|,
 \qquad 0\le\bm r\le\bm n,
 \quad |\bm r|=r_\up+r_\dn.
\end{equation}
For a fixed total number \(N\le(2-\epsilon)M\), one has
\begin{equation}\label{eq:removaltotal}
 E_{\mathrm{tot}}(N-r)
       \ge E_{\mathrm{tot}}(N)-C_\epsilon r,
 \qquad 0\le r\le N.
\end{equation}
\end{lemma}
\begin{proof}
Add particles one at a time from \(\bm n-\bm r\) to \(\bm n\).
At every step in \eqref{eq:addspin}, the denominator is at least
\(\epsilon M\).  Summing the upper bounds on the addition energies
and rearranging proves \eqref{eq:removalspin}.
Likewise, along the sequence \(N-r,\ldots,N\), the denominators
in \eqref{eq:addtotal} are at least \(\epsilon M\), giving
\eqref{eq:removaltotal}.  These are one-sided energy estimates;
they require neither positivity nor gaps in any intermediate sector.
\end{proof}

\begin{lemma}[General spectral-band comparison]\label{lem:generalcomparison}
Let \(H\) satisfy Definition~\ref{def:local}.  Suppose its lowest
\(d\) states in the selected all-LLL sector have width
\(E_d-E_1\le w\), with a gap \(E_{d+1}-E_d\ge\gamma>0\)
above them.  Write \(P\) for their spectral projection.  For this
comparison alone, no relation between \(\gamma\) and \(w\) is needed.
Assume a size-independent constant \(C_{\mathrm{rem}}\ge0\)
bounds the decrease of ground energy on removing particles:
\eqref{eq:removalspin} in a fixed-spin sector, or
\eqref{eq:removaltotal} in a fixed-total-number sector, with
\(C_\epsilon\) replaced by \(C_{\mathrm{rem}}\).
Choose a fixed \(\mu>C_{\mathrm{rem}}+w\) in \eqref{eq:penalty}.
Then the enlarged Hamiltonian has exactly the same lowest \(d\)
states, with the same energies and projection \(P\), and the gap
above this band is at least
\begin{equation}\label{eq:generalgap}
 \Delta=\min\{\gamma,\mu-C_{\mathrm{rem}}-w\}>0.
\end{equation}
Its spectrum strictly below \(E_1+\mu-C_{\mathrm{rem}}\),
including multiplicities, is the corresponding all-LLL spectrum.
In particular, taking \(\mu>C_{\mathrm{rem}}+w+\gamma\)
retains the lower bound \(\Delta=\gamma\) on the band gap.
\end{lemma}
\begin{proof}
At fixed spin populations \(\bm n\), the block with
\(\bm r\) spectator particles has interaction energies bounded
below by \(E_0(\bm n-\bm r)\), and the penalty is
\(\mu|\bm r|\).  Since \(E_0(\bm n)=E_1\), its bottom obeys
\[
 E_0(\bm n-\bm r)+\mu|\bm r|
       \ge E_1+(\mu-C_{\mathrm{rem}})|\bm r|.
\]
Every block with at least one spectator lies above
\(E_1+\mu-C_{\mathrm{rem}}\), or at that energy.  The block
with no spectators is exactly the original LLL Hamiltonian.
Thus all selected \(d\) levels remain below every spectator block.
Subtracting their top energy \(E_d\le E_1+w\) gives a spectator
gap of at least \(\mu-C_{\mathrm{rem}}-w\); the remaining all-LLL
states lie at least \(\gamma\) above \(E_d\).
This proves the assertions in the fixed-spin case.
At fixed total number \(N\), a block with \(r\) spectators
has interaction energy at least \(E_{\mathrm{tot}}(N-r)\);
use \(E_{\mathrm{tot}}(N)=E_1\) and \eqref{eq:removaltotal}
to obtain the same comparison.
\end{proof}

Lemma~\ref{lem:removal} verifies this comparison for general local
polynomials below saturation, with \(C_{\mathrm{rem}}=C_\epsilon\).
For a positive parent with \(E_1=0\), positivity in every number
sector gives the removal bound directly with
\(C_{\mathrm{rem}}=0\), recovering Lemma~\ref{lem:comparison}
without a filling restriction.  The particle-addition argument
does not establish a uniform constant when the density of empty
orbitals tends to zero.  The positive-parent comparison remains
available in that case, or one can verify a uniform removal bound
by another argument.

For balanced populations with limiting filling \(0<\nu<2\),
one can choose any fixed \(\epsilon<1-\nu/2\) for all large
sizes.  For the full fixed-total-number problem it suffices to
choose \(\epsilon<2-\nu\); individual spin sectors may then be
saturated.  A scalar shift of the Hamiltonian does not replace
this energy comparison: it shifts all the number sectors by the
same constant and leaves their energy differences unchanged.

For neutral observables the dynamics of \(H_{\mathrm{ext}}\)
is the dynamics of the bounded ambient operator
\begin{equation}\label{eq:Hb}
 H_b=H_{\mathrm{LLL}}-\mu\sum_\sigma\dG_\sigma(\Pi)
     =\int_{\TT_L}k_R\dd R,\qquad
 k_R=h_R-\mu\sum_\sigma
           a_\sigma^\dagger(e_R)a_\sigma(e_R).
\end{equation}
Indeed, \(H_{\mathrm{ext}}=H_b+\mu(N_\up+N_\dn)\).
Every \(k_R\) is bounded, neutral, and has Gaussian local
approximants, by \eqref{eq:tails}.  No estimate involves a point
annihilator on the full continuum space.
In the selected sector of total number \(N\), its low-energy band
lies in \([E_1-\mu N,E_d-\mu N]\), with unchanged width and gap
above it.  The filters below use energy differences only.

\begin{remark}[Why the physical charge is retained]
On the LLL, the compression of \(Q_X\) is
\(\dG_\up(\Pi1_X\Pi)\).  In general \((\Pi1_X\Pi)^2\ne\Pi1_X\Pi\),
so its spectrum need not be integral.  We do not use this compressed
charge in the flux exponentials.  The physical charge \(Q_X\)
can take a low-energy vector out of the LLL; the penalty makes
these excursions compatible with a gap in the enlarged sector.
\end{remark}

\subsection{Verification from an isolated band at the original boundary condition}

\begin{proposition}[The analytical input]\label{prop:analytic}
Let \(H_{\mathrm{LLL}}\) satisfy Definition~\ref{def:local}.
Suppose its lowest \(d\) states in the all-LLL sector
\((n_\up,n_\dn)\) satisfy Assumption~\ref{ass:band}, with
width bound \(w\), gap bound \(\gamma>2w\), and projection \(P\).
Assume either
\begin{enumerate}
\item it is a positive parent of Definition~\ref{def:parent}
and \(E_1=0\); or
\item \(n_\sigma\le(1-\epsilon)M\) for both spins, with
a fixed \(\epsilon>0\).
\end{enumerate}
More generally, any uniform removal bound as in
Lemma~\ref{lem:generalcomparison} suffices.
Choose \(\mu>C_{\mathrm{rem}}+w+\gamma\), with \(C_{\mathrm{rem}}=0\)
in the first case and \(C_{\mathrm{rem}}=C_\epsilon\) in the
second.  Then:
\begin{enumerate}
\item The ambient dynamics generated by \(H_b\) has a
size-independent Lieb--Robinson bound for bounded local even
observables.
\item For \(Q=Q_\Gamma\), the current is bounded on ambient Fock
space and splits as
\[
 \ii[H_b,Q]=J_-+J_+,
\]
where \(J_\pm\) have norm \(O(L)\) and Gaussian local approximants
near \(\partial_\pm\).
\item There are bounded, neutral boundary dressings \(K_\pm\),
supported in strips of radius \(L_x/64\), with polynomial ambient
norms, such that
\[
 [Q-K_--K_+,P]\eqL0.
\]
\item The low-energy-sector clustering estimate
\eqref{eq:groundcluster} holds for bounded neutral local operators.
\end{enumerate}
\end{proposition}

\begin{proof}
The proofs are given in Appendices~\ref{app:locality} and
\ref{app:filters}.  We spell out the role of the sector here.
Local approximants to \(k_R\), and to the currents, preserve both
particle numbers.  A spectral filter applied to these operators
therefore connects states only inside \(\HH_{\bm n}\).  The gap
\(\Delta=\gamma>2w\) above this band, retained by
Lemma~\ref{lem:generalcomparison}, is sufficient, even though
no grand-canonical gap is assumed.  The commutator
and propagation estimates use the polynomial bounds in
Definition~\ref{def:local}, with no positivity requirement.
The inverse-energy filter handles transitions between the band and
its complement.  The clustering filter is zero on all intraband
energy differences, which have magnitude at most \(w<\Delta/2\).
The projection identities are therefore exact even when the
levels inside \(P\) are split; neither equality of their energies
nor a vanishing band width is used.
Ambient norm estimates supply
locality, while the charge estimates
\(\norm{Q_X|_{\HH_{\bm n}}}\le CL^2\) control the algebraic
errors involving charges.  This division is used throughout.
\end{proof}

For a unitary moving physical supports by a uniformly bounded
distance, such as the magnetic translation below, conjugating
these \(K_\pm\) gives the \(K_\pm^U\) required in
Assumption~\ref{ass:index}.  A unitary with rapidly decaying
locality tails works as well if its neutral local approximants
have uniform polynomial bounds.

\section{Examples and the filling constraint}

\subsection{Magnetic translation}

Let \(\delta=L_x/M=2\pi/L_y\), and define the one-particle magnetic
translation
\[
 (t_xf)(x,y)=e^{\ii\delta y}f(x-\delta,y).
\]
It respects \eqref{eq:bc}, commutes with \(\Pi\), and moves
spatial supports by \(\delta\).  Write
\(\Gamma_{\mathrm a}(v)=\bigoplus_{k\ge0}\wedge^k v\) for
fermionic second quantization, distinct from the region \(\Gamma\).
Take
\(U=\Gamma_{\mathrm a}(t_x\otimes1_{\mathbb C^2})^\dagger\), so that
\begin{equation}\label{eq:shiftQ}
 U^\dagger Q_XU=Q_{X+\delta\hat x}.
\end{equation}
Assume the interaction is invariant under this common
translation of both spins.  Then
\([U,H_b]=[U,H_{\mathrm{ext}}]=[U,P]=0\).

For \(M>2\), let
\[
 S_-^0=[0,\delta)\times S_y^1,\qquad
 S_+^0=[L_x/2,L_x/2+\delta)\times S_y^1.
\]
Equation~\eqref{eq:shiftQ} gives exactly
\[
 U^\dagger QU-Q=T_-+T_+,\qquad
 T_-=-Q_{S_-^0},\quad T_+=Q_{S_+^0}.
\]
For large \(L\) these strips lie in \(B_\pm\).  Also
\(Q_-+T_-\) counts the region obtained by removing \(S_-^0\)
from \(\Gamma\cap S_-\), and \(Q_++T_+\) counts the region
obtained by adding \(S_+^0\) to \(\Gamma\cap S_+\).  Thus
\eqref{eq:transportnormalization} is exact.  Their sector
norms are at most \(n_\up\).

\begin{theorem}[LLL filling constraint]\label{thm:filling}
Let the Hamiltonians, selected spin-number sectors, and low-energy
band projections satisfy Proposition~\ref{prop:analytic}, with fixed
band rank \(d\), and assume invariance under common
magnetic translation.  Then, for \(\sigma=\up,\dn\),
\begin{equation}\label{eq:filling}
 \dist\left(\frac{dn_\sigma}{M},\mathbb Z\right)
       =O(L^{-\infty}).
\end{equation}
No spin-rotation symmetry or equality of the two populations is
needed for this statement.
For all sufficiently large tori these constraints imply the exact
divisibility
\begin{equation}\label{eq:finiteinteger}
 M\mid dn_\up,\qquad M\mid dn_\dn,
 \qquad\text{equivalently}\qquad
 \frac{M}{\gcd(M,n_\up,n_\dn)}\mid d.
\end{equation}
\end{theorem}
\begin{proof}
Proposition~\ref{prop:analytic} and the exact transport calculation
verify Assumption~\ref{ass:index}.  The \(M\) translates of
\(S_-^0\) partition the torus.  Since \(UPU^\dagger=P\),
their charge traces are equal, and hence
\[
 M\Tr(PQ_{S_-^0})=\Tr(PN_\up)=dn_\up.
\]
Insert \(\Tr(PT_-)=-dn_\up/M\) into Theorem~\ref{thm:index}.
Replacing the spin-up charge throughout by the spin-down charge
proves the second assertion.
For the final statement, a noninteger \(dn_\sigma/M\) is at
distance at least \(1/M\) from the integers.  Since \(M=O(L^2)\),
this is incompatible with \eqref{eq:filling} for large \(L\).
The equivalent greatest-common-divisor formulation follows from
B\'ezout's identity applied to \(M,n_\up,n_\dn\).
\end{proof}

\subsection{Rational fillings and the parity distinction}

\begin{corollary}[Separate spin fillings]\label{cor:separate}
Under the hypotheses of Theorem~\ref{thm:filling}, suppose
\[
 \frac{n_\sigma}{M}\longrightarrow\nu_\sigma
       =\frac{a_\sigma}{b_\sigma},\qquad
 \gcd(a_\sigma,b_\sigma)=1,\quad b_\sigma>0.
\]
Then
\[
 \operatorname{lcm}(b_\up,b_\dn)\mid d.
\]
In particular, the conclusion holds when the two rational
fillings are exactly fixed at every size.
\end{corollary}
\begin{proof}
Taking the limit of \eqref{eq:filling}, the fixed number
\(d\nu_\sigma\) is an integer.  Coprimality gives
\(b_\sigma\mid d\) for each spin, and hence the asserted
least common multiple.  The two constraints do not in general
multiply: when the spin fillings coincide, they are the same
arithmetic condition.
\end{proof}

\begin{corollary}[Balanced rational filling]\label{cor:rational}
Under the hypotheses of Theorem~\ref{thm:filling}, suppose
\[
 n_\up=n_\dn=N/2,\qquad
 \frac{N}{M}\longrightarrow\nu=\frac pq,\qquad
 0<\nu\le2,\quad p,q\in\mathbb N,\quad\gcd(p,q)=1.
\]
Then the low-energy band rank in this balanced sector obeys
\begin{equation}\label{eq:generaldivisor}
 \boxed{\quad b_\nu\mid d,\qquad
 b_\nu:=\frac{2q}{\gcd(p,2q)}
       =\frac{2q}{\gcd(p,2)}
       =\begin{cases}
          2q,&p\text{ odd},\\
          q,&p\text{ even}.
         \end{cases}\quad}
\end{equation}
\end{corollary}
\begin{proof}
The spin filling is \(\nu_\sigma=p/(2q)\).  Applying
Corollary~\ref{cor:separate}, or taking the limit directly in
\eqref{eq:filling}, gives
\begin{equation}\label{eq:balancedinteger}
 \frac{dp}{2q}\in\mathbb Z.
\end{equation}
The denominator of \(p/(2q)\) in lowest terms is
\(2q/\gcd(p,2q)\).  Since \(p\) and \(q\) are coprime,
\(\gcd(p,2q)=\gcd(p,2)\), proving the result.
\end{proof}

For \(0<\nu<2\), the balanced-population hypothesis makes the
partial-filling condition in Proposition~\ref{prop:analytic}
automatic for all large sizes.  Thus this corollary applies to
general local LLL field polynomials with an isolated low-energy
band satisfying Assumption~\ref{ass:band}, at any lowest energy
and with possible splitting inside the band.  At the endpoint \(\nu=2\),
the positive-parent comparison or another uniform removal bound
is the available hypothesis for this argument.

The parity distinction can equivalently be organized by the
denominator of the total filling:
\[
\begin{array}{c|c|c}
 q & p\text{ (coprime to }q\text{)} & \text{divisor forced on }d\\ \hline
 \text{even}&\text{necessarily odd}&2q\\
 \text{odd}&\text{odd}&2q\\
 \text{odd}&\text{even}&q
\end{array}
\]
Thus even \(q\) always gives a multiple of \(2q\).  For odd
\(q\), the parity of \(p\) decides whether the factor of two
survives.  Section~\ref{sec:exacttranslation} compares examples
with the exact translation degeneracy.  These are necessary
divisors; the theorem does not assert the existence of a gapped
state attaining any of these dimensions.

There is no second factor of two from applying the theorem to
spin down: at balance it reproduces \eqref{eq:balancedinteger}.
Using total charge gives \(dp/q\in\mathbb Z\), which is already
implied by that equation.  More generally, integer combinations
of the two spin charges give integer combinations of the same
two transport indices and do not strengthen this arithmetic.

\begin{remark}[Admissible finite tori]\label{rem:finiteparity}
If \(N/M=p/q\) is imposed exactly and the populations are
balanced, the allowed sizes are
\[
 M=b_\nu k,\qquad
 n_\up=n_\dn=\frac{p}{\gcd(p,2)}k,
 \qquad k\in\mathbb N.
\]
Indeed, integrality of \(n_\sigma=Mp/(2q)\) is equivalent to
\(b_\nu\mid M\).  In particular, when \(p\) is odd one must
take an even multiple of \(q\) flux quanta to have even \(N\).
This finite-size restriction does not supply another factor in
the low-energy band rank.  The limiting-filling formulation
does not assume exact rational filling at each finite size.
Under the band hypotheses, \eqref{eq:finiteinteger} implies that each
spin filling belongs eventually to the finite set
\(\{0,1/d,\ldots,1\}\); convergence then forces it to equal its
limit for all sufficiently large sizes.
\end{remark}

\subsection{Comparison with exact translation degeneracy}
\label{sec:exacttranslation}

For this comparison, impose symmetry also under the common
\(y\)-translation.  The additional one-particle unitary is
\[
 (t_yf)(x,y)=f(x,y-\eta),\qquad \eta=L_y/M.
\]
It respects \eqref{eq:bc} because \(L_x\eta=2\pi\), and it
commutes with the covariant derivatives and \(\Pi\).
Direct substitution gives
\[
 t_xt_y=e^{\ii\delta\eta}t_yt_x
        =e^{2\pi\ii/M}t_yt_x.
\]
Set \(\mathcal T_\alpha=\Gamma_{\mathrm a}
(t_\alpha\otimes1_{\mathbb C^2})\), for \(\alpha=x,y\).
Thus \(\mathcal T_x=U^\dagger\) in the convention used above.
On a fixed-total-number sector, and also on any fixed-spin-number
sector with total number \(N\), their algebra is
\begin{equation}\label{eq:COMalgebra}
 \mathcal T_x\mathcal T_y
       =e^{2\pi\ii N/M}\mathcal T_y\mathcal T_x.
\end{equation}

\begin{proposition}[Exact center-of-mass degeneracy]
\label{prop:COM}
If an LLL Hamiltonian commutes with both \(\mathcal T_x\) and
\(\mathcal T_y\), the dimension \(d_E\) of each energy eigenspace
in either of these number sectors satisfies
\[
 q_L:=\frac{M}{\gcd(M,N)}\mid d_E.
\]
This is a finite-size symmetry statement, requiring neither a
gap nor locality nor positivity of the interaction.
\end{proposition}
\begin{proof}
The eigenspace is finite-dimensional and invariant under the two
translations.  Take the determinant of \eqref{eq:COMalgebra}
restricted to it.  Cancelling the determinants of the unitaries
gives \(e^{2\pi\ii Nd_E/M}=1\), or \(M\mid Nd_E\).
This is equivalent to the stated divisibility.
\end{proof}

This is the center-of-mass part of Haldane's magnetic-translation
analysis \cite{Haldane85}.  An isolated band contains complete
eigenspaces, so its rank also obeys \(q_L\mid d\).
At exact total filling \(N/M=p/q\) in lowest terms, \(q_L=q\).
The main index theorem requires
only the common \(x\)-translation; when both translations are
symmetries, its improvement can be quantified.  For an isolated
low-energy band as in Assumption~\ref{ass:band} in a fixed spin-number sector,
\eqref{eq:finiteinteger} gives, for all sufficiently large sizes,
\[
 b_L:=\frac{M}{\gcd(M,n_\up,n_\dn)}\mid d,
 \qquad
 \frac{b_L}{q_L}
    =\frac{\gcd(M,N)}{\gcd(M,n_\up,n_\dn)}.
\]
The quotient is an integer because a common divisor of
\(M,n_\up,n_\dn\) also divides \(N=n_\up+n_\dn\).
At balanced exact filling, \(b_L=b_\nu\), so the gain is
\(2/\gcd(p,2)\): a factor of two for odd \(p\), and no gain
in this divisibility constraint for even \(p\).
For a fully polarized sector, say \(n_\dn=0\), one has
\(b_L=q_L\) as well.
\[
\begin{array}{c|c|c}
 \nu\text{ (balanced)}
 &\text{exact symmetry divisor }q
 &\text{gapped index divisor }b_\nu\\ \hline
 1/2&2&4\\
 1/3&3&6\\
 2/3&3&3\\
 2/5&5&5\\
 3/4&4&8\\
 1&1&2
\end{array}
\]
The proof of Proposition~\ref{prop:COM} is included here so that
this comparison does not require a result from the cited paper.
The two conclusions concern different spectral objects: the
translation algebra constrains each exact energy eigenspace,
whereas the index constrains the whole isolated band.  For
example, at \(\nu=1/2\), a rank-four low-energy band split into
two translation doublets is compatible with both conclusions.

\subsection{The full low-energy band with spin-exchange symmetry}

Specifying spin-\(1/2\) particles alone does not imply balanced
populations.  One may impose the balanced sector as above.
Alternatively, a spin symmetry gives precisely the trace
identity needed for the full low-energy band, without requiring
every vector in it to be a singlet.

\begin{theorem}[Spin-symmetric fixed-total-number version]
\label{thm:spinexchange}
Let the Hamiltonians of Definition~\ref{def:local} be invariant
under common magnetic translation.  In the full all-LLL sector
of fixed total particle number \(N=N_\up+N_\dn\), suppose
the lowest \(d\) states satisfy Assumption~\ref{ass:band}:
\(d\) is fixed, their width is at most \(w\), and the gap above
them is at least \(\gamma>2w\).  Let \(P\) be the spectral
projection onto this band, including all its spin sectors.
Assume also either that the Hamiltonian is
a positive parent with \(E_1=0\), or that
\(N\le(2-\epsilon)M\) for a fixed \(\epsilon>0\).
More generally, a uniform bound of the form
\eqref{eq:removaltotal} suffices in place of these alternatives.
Assume an ambient unitary \(R\) exchanges the two spin numbers
and preserves the penalty:
\[
 [R,H_{\mathrm{LLL}}]=0,\qquad
 \left[R,\sum_\sigma\dG_\sigma(1-\Pi)\right]=0,
\]
\[
 R^\dagger N_\up R=N_\dn,\qquad
 R^\dagger N_\dn R=N_\up.
\]
Then
\begin{equation}\label{eq:spintracefilling}
 \dist\left(\frac{dN}{2M},\mathbb Z\right)=O(L^{-\infty}).
\end{equation}
In particular, \(2M\mid dN\) for all sufficiently large tori.
If \(N/M\to p/q\) in lowest terms, the same divisor
\(b_\nu=2q/\gcd(p,2)\) in \eqref{eq:generaldivisor} divides
this full low-energy band rank.  Charge conservation and
\(\mathrm{SU}(2)\) spin symmetry imply the spin conservation
and exchange assumptions used here.
\end{theorem}
\begin{proof}
The index proof and its analytical input apply on
\[
 \HH_N=\{\psi:(N_\up+N_\dn)\psi=N\psi\}
\]
in place of \(\HH_{\bm n}\).  To verify this extension, note
that all local interaction terms, currents, dressings, and
unitaries used in the proof are neutral and thus preserve
\(\HH_N\).  The physical charges obey \(0\le Q_X\le N\)
on this sector, so every polynomial norm estimate still holds.
Use the total-number version of Lemma~\ref{lem:generalcomparison}
and choose \(\mu>C_{\mathrm{rem}}+w+\gamma\).
The enlarged Hamiltonian has the same lowest \(d\) states,
with their original energies and projection \(P\), and a gap
at least \(\Delta=\gamma>2w\) above them in \(\HH_N\).
This gives the same spectral-filter identities and clustering
estimates, including when the band is split.  In particular,
the filter does not need a gap to states in different total-number
sectors.  In the partial-filling case, the energy comparison uses
only the variational bounds of Lemma~\ref{lem:removal} in those
other sectors; the positive-parent case uses their nonnegative
energies directly.

The only step of Theorem~\ref{thm:filling} that used a definite
spin number was its final trace evaluation.  The strip partition
now gives, instead,
\[
 M\Tr(PQ_{S_-^0})=\Tr(PN_\up),\qquad
 \Tr(PT_-)=-\frac1M\Tr(PN_\up).
\]
The index theorem makes this last quantity an integer up to
\(O(L^{-\infty})\).  The two commutation assumptions on \(R\)
imply \([R,H_{\mathrm{ext}}]=0\), and spin exchange preserves
the fixed-total-number sector.  The gap comparison identifies
\(P\) with the spectral projection onto the whole low-energy
band of \(H_{\mathrm{ext}}\) in that sector; hence
\(R^\dagger PR=P\).
Cyclicity and the exchange symmetry give
\[
 \Tr(PN_\up)=\Tr(PN_\dn),\qquad
 \Tr\bigl(P(N_\up+N_\dn)\bigr)=dN.
\]
Consequently each spin trace is \(dN/2\), proving
\eqref{eq:spintracefilling}.  Its left side is either zero or
at least \(1/(2M)\), which proves the exact eventual divisibility.
The same arithmetic as in
Corollary~\ref{cor:rational} proves the divisibility statement.
For an \(\mathrm{SU}(2)\)-invariant Hamiltonian, the second
quantization of a rotation by \(\pi\) about the spin \(x\)-axis
supplies \(R\).  It acts only on spin, so it commutes with the
spin-independent penalty.  Conservation of \(N\) and \(S^z\)
gives conservation of both spin numbers.
\end{proof}

For this theorem the statement \(\nu_\up=\nu_\dn=\nu/2\)
means the filling in the normalized band trace:
\[
 \frac{\Tr(PN_\sigma)}{dM}=\frac{N}{2M}.
\]
It need not hold as a sharp particle number in each vector
in the band.  No evenness assumption on \(N\) is needed here;
if exact filling is imposed, \(M=qk\), \(N=pk\) suffices.
By contrast, Corollary~\ref{cor:rational} concerns the balanced
sector alone and requires a gap only within that sector.  The
two ranks agree when the same low-energy interval is used and
every state in the full fixed-\(N\) band is a singlet, but should
otherwise be distinguished.

\subsection{Special cases and the gap obstruction}

\begin{corollary}[Balanced half filling]\label{cor:four}
At \(\nu=1/2\), either the balanced-sector hypotheses of
Corollary~\ref{cor:rational} or the full-sector hypotheses of
Theorem~\ref{thm:spinexchange} imply \(4\mid d\).
\end{corollary}
\begin{proof}
Here \(p=1\), \(q=2\), and hence \(b_\nu=4\).
\end{proof}

\begin{corollary}[Obstruction for a collapsing low-energy band]
\label{cor:gapobstruction}
Fix a rational limiting filling \(\nu=p/q\) and either the
balanced-sector or the full spin-exchange-symmetric setting.
Retain the locality, symmetry, and uniform energy-comparison
hypotheses, but do not assume a uniform spectral gap.
Suppose for a fixed \(d\) that
\[
 w_L=E_d(L)-E_1(L)\longrightarrow0,\qquad b_\nu\nmid d.
\]
Then
\begin{equation}\label{eq:gapobstruction}
 E_{d+1}(L)-E_d(L)\longrightarrow0,
 \qquad E_{d+1}(L)-E_1(L)\longrightarrow0.
\end{equation}
\end{corollary}
\begin{proof}
If the first limit failed, some subsequence would have
\(E_{d+1}-E_d\ge\gamma_0>0\).  Since \(w_L\to0\),
choose a fixed \(w<\gamma_0/2\) bounding \(w_L\) eventually
along that subsequence.  Assumption~\ref{ass:band} then holds
there, and the corresponding divisibility theorem gives
\(b_\nu\mid d\), a contradiction.  Adding \(w_L\to0\)
gives the second limit.
\end{proof}

For a positive parent at balanced half filling, exactly ten
zero modes give \(d=10\), \(w_L=0\), and hence \(E_{11}(L)\to0\).
The same obstruction applies to a proposed rank-ten band with
exponentially small, or merely vanishing, splitting.  Spectral
counting is a separate input; no assertion about that count is
used or proved in this note.

It is essential to distinguish these two uses of \(d\).
Ten exact zero modes alone do not exclude a larger isolated
low-energy band: additional positive-energy states could tend
to zero and produce a band of rank twelve, sixteen, or another
multiple of four.  The conclusion from the exact zero-mode
count alone is \(E_{11}\to0\).  To exclude a uniformly gapped
phase with a finite quasi-degenerate low-energy sector by this
argument, one must identify the rank of that entire sector and
show that it violates the divisor, or otherwise exclude all
such larger bands.

When both common magnetic translations are symmetries, the exact
algebra at \(\nu=1/2\) requires only \(2\mid d\), which is
compatible with ten zero modes.  The additional band constraint
\(4\mid d\) contradicts a uniform gap above a rank-ten band
whose width vanishes.  This identifies both the gap assumption
used for contradiction and the spectral sector to which it applies.

The distinction for even \(p\) is necessary even for a
spin-singlet ground state.  For example, at \(\nu=2\) the
two spin components of the LLL are completely filled.  The
all-LLL sector with \(N=2M\) is one-dimensional and its filled
state is a spin singlet.  Take the allowed positive parent
\(H_{\mathrm{LLL}}=0\); after the enlargement its Hamiltonian
is \(\mu\sum_\sigma\dG_\sigma(1-\Pi)\), with this unique
ground state and gap \(\mu\) at fixed \(N\).  The all-LLL
sector itself has no orthogonal states.  Thus a universal
claim \(2q\mid d\) would already fail at \(p=2,q=1,d=1\),
whereas \eqref{eq:generaldivisor} gives the correct divisor one.

\subsection{A spin-triplet derivative parent}

A concrete parent in the class is specified by
\[
 \psi_\sigma(R)=a_\sigma(e_R),\qquad
 \eta_\sigma(R)=a_\sigma(g_R),
\]
\begin{align*}
 B_+(R)&=\psi_\up(R)\eta_\up(R),\\
 B_-(R)&=\psi_\dn(R)\eta_\dn(R),\\
 B_0(R)&=2^{-1/2}\bigl(
       \psi_\dn(R)\eta_\up(R)+\psi_\up(R)\eta_\dn(R)\bigr),\\
 H_{\mathrm{trip}}&=8\pi\int_{\TT_L}
       \bigl(B_+^\dagger B_++B_0^\dagger B_0+B_-^\dagger B_-\bigr)\dd R.
\end{align*}
The positive normalization is immaterial for the index.  Each
summand preserves the two spin numbers and is covariant under
common magnetic translation.  The coefficients select a
spin-triplet pair and its first relative derivative at coincidence.
All the hypotheses on the interaction follow directly from
\eqref{eq:tails}; a gap above the proposed low-energy sector
is assumed only for the contradiction argument.

\section{Proof of the index theorem}\label{sec:proof}

We work on \(\HH_{\bm n}\) when manipulating charges, and continue
to use ambient supports for local unitaries.  On this sector all
charges and their bounded perturbations have polynomial norms.
The proof uses the rank and approximate invariance of \(P\),
the dressing identity, and clustering.  It never requires the
Hamiltonian to be a scalar on \(\Ran P\).
Every error below is uniform for \(0\le\phi\le2\pi\).

\subsection{Approximate invariant subspaces}

\begin{lemma}[Separated factors]\label{lem:factors}
Let \(V_\pm\) be neutral ambient unitaries supported in \(S_\pm\).
If \([V_-V_+,P]\eqL0\) and \eqref{eq:groundcluster} holds, then
\([V_\pm,P]\eqL0\).
\end{lemma}
\begin{proof}
Write \(V=V_+V_-\); the factors commute.  For every unit vector
\(\psi\in\Ran P\), approximate invariance gives
\(\norm{PV\psi}\ge1-\epsL\).  Clustering gives
\[
 PV_+V_-P\eqL PV_+PV_-P.
\]
Here the order of supports is reversed relative to
\eqref{eq:groundcluster}; the required estimate follows from
\[
 \norm{PV_+(1-P)V_-P}
 =\norm{PV_-^\dagger(1-P)V_+^\dagger P}=O(L^{-\infty}).
\]
Since \(\norm{PV_+P}\le1\), it follows that
\(\norm{PV_-\psi}\ge1-\epsL\).
Consequently \(\norm{(1-P)V_-P}=\epsL\).
Apply the same argument to \(V^\dagger\), reordering its
separated factors, to get
\(\norm{(1-P)V_-^\dagger P}=\epsL\).
These are the two off-diagonal blocks of \([V_-,P]\).
Interchanging the two signs proves the other assertion.
\end{proof}

\begin{lemma}[Compression of an almost invariant unitary]\label{lem:polar}
If \(V\) is unitary and \([V,P]\eqL0\), then
\(PVP|_{\Ran P}\) differs by \(\epsL\) from a unitary on
\(\Ran P\).  If \(A=A^\dagger\) is sector-bounded and
\([A,P]\eqL0\), then \([e^{\ii\phi A},P]\eqL0\).
\end{lemma}
\begin{proof}
For \(v=PVP|_{\Ran P}\),
\[
 v^\dagger v=P-PV^\dagger(1-P)VP\eqL P.
\]
Its polar factor \(v(v^\dagger v)^{-1/2}\) gives the first claim.
For the second use
\[
 [e^{\ii\phi A},P]
 =\ii\int_0^\phi e^{\ii sA}[A,P]e^{\ii(\phi-s)A}\dd s.
\]
The unitaries have norm one, so the estimate does not grow
exponentially with \(\norm A\).
\end{proof}

\subsection{Local charge fluctuations}

Define
\[
 \bar Q_\pm=Q_\pm-K_\pm,\quad
 \bar Q=Q-K_--K_+=\bar Q_-+Q_m+\bar Q_+,
\]
and
\[
 \bar Q_\pm^U=Q_\pm+T_\pm-K_\pm^U,\qquad
 \bar Q^U=\bar Q_-^U+Q_m+\bar Q_+^U.
\]
The support choices ensure that objects with opposite boundary
signs commute and that \(Q_m\) commutes with all boundary objects.
In particular,
\begin{equation}\label{eq:dressedfacts}
 \bar Q^U\eqL U^\dagger\bar Q U,\qquad
 [\bar Q,P]\eqL0,\qquad [\bar Q^U,P]\eqL0.
\end{equation}
All of these statements are sector norm statements.  Exponentials
of the local operators also have ambient local realizations:
their unbounded parts are the local charges in
\eqref{eq:transportnormalization}, and their perturbations are bounded.

Set
\begin{equation}\label{eq:Z}
 Z(\phi)=e^{\ii\phi\bar Q^U}e^{-\ii\phi\bar Q},\qquad
 Z_\pm(\phi)=e^{\ii\phi\bar Q_\pm^U}e^{-\ii\phi\bar Q_\pm}.
\end{equation}
The commuting spatial decomposition gives \(Z=Z_-Z_+\).
Equation~\eqref{eq:dressedfacts} gives \([Z,P]\eqL0\).
Lemma~\ref{lem:factors} therefore implies
\begin{equation}\label{eq:Zpreserve}
 [Z_\pm(\phi),P]\eqL0.
\end{equation}
Similarly, since \(Q_m\) is integer-valued,
\[
 e^{2\pi\ii\bar Q}
   =e^{2\pi\ii\bar Q_-}e^{2\pi\ii\bar Q_+}.
\]
Another application of Lemma~\ref{lem:factors} yields
\begin{equation}\label{eq:Vpreserve}
 [e^{2\pi\ii\bar Q_\pm},P]\eqL0.
\end{equation}

\subsection{A formula for the determinant}

On the \(d\)-dimensional space \(\Ran P\), set
\[
 \mathcal Z_-(\phi)=PZ_-(\phi)P,\qquad
 D_-=\bar Q_-^U-\bar Q_-=T_--K_-^U+K_-.
\]
By \eqref{eq:Zpreserve}, \(\mathcal Z_-\) is uniformly close to
a unitary and hence invertible, with uniformly bounded inverse.
Differentiating the bounded sector operators gives
\begin{align}
 -\ii\mathcal Z_-'(\phi)
 &=PZ_-(\phi)e^{\ii\phi\bar Q_-}
                D_-e^{-\ii\phi\bar Q_-}P \notag\\
 &\eqL\mathcal Z_-(\phi)
       P e^{\ii\phi\bar Q}D_-e^{-\ii\phi\bar Q}P.
       \label{eq:Zderivative}
\end{align}
In the second line one first inserts \(P\) after \(Z_-\), using
\eqref{eq:Zpreserve}; the remaining parts of \(\bar Q\) commute
with \(D_-\).  The polynomial norm of \(D_-\) does not change
the error class.  Since \([\bar Q,P]\eqL0\), the trace of the
last factor is \(\Tr(PD_-)+\epsL\).

Jacobi's formula for the determinant now gives
\[
 \frac{\dd}{\dd\phi}\det\mathcal Z_-(\phi)
   =\bigl(\ii\Tr(PD_-)+\epsL\bigr)
                      \det\mathcal Z_-(\phi).
\]
Because \(D_-\) is self-adjoint and \(\mathcal Z_-(0)=1_{\Ran P}\),
integration over a bounded interval yields
\begin{equation}\label{eq:detformula}
 \det\mathcal Z_-(\phi)
       =e^{\ii\phi\Tr(PD_-)}+\epsL
       =e^{\ii\phi\Tr(PT_-)}+\epsL.
\end{equation}
For the last equality use \(K_-^U\eqL U^\dagger K_-U\) and
\([U,P]\eqL0\).  Cyclicity of the finite-rank trace gives
\(\Tr(PK_-^U)-\Tr(PK_-)=\epsL\).
This is the finite-rank counterpart of computing the phase
acquired by a single invariant state.

\subsection{The endpoint and integrality}

Let \(V=e^{2\pi\ii\bar Q_-}\).  The integer spectra of \(Q_m,Q_+\)
give
\[
 V=e^{2\pi\ii(Q-K_-)}.
\]
Conjugate by \(U\), use the transport identity and
\(K_-^U\eqL U^\dagger K_-U\), and split the commuting boundary
terms.  The normalization
\(e^{2\pi\ii(Q_++T_+)}=1\) gives
\[
 U^\dagger VU
 \eqL e^{2\pi\ii(Q+T_-+T_+-K_-^U)}
 =e^{2\pi\ii\bar Q_-^U}.
\]
It follows that
\[
 Z_-(2\pi)\eqL U^\dagger VU V^\dagger.
\]
Both \(U\) and \(V\) preserve \(P\) up to \(\epsL\), by
Assumption~\ref{ass:index} and \eqref{eq:Vpreserve}.
After compression and replacement by their polar unitaries
(Lemma~\ref{lem:polar}), the right-hand side is a commutator
of two \(d\)-dimensional unitaries, up to \(\epsL\).
Its determinant is therefore \(1+\epsL\).
Equation~\eqref{eq:detformula} implies
\[
 e^{2\pi\ii\Tr(PT_-)}=1+\epsL.
\]
For real \(x\), distance from \(x\) to \(\mathbb Z\) is bounded
by a constant times \(\abs{e^{2\pi\ii x}-1}\).
This proves Theorem~\ref{thm:index}.

\appendix
\section{Gaussian LLL kernels on the torus}\label{app:kernel}

Let \(\kappa=2\pi/L_y\).  An orthonormal basis of the LLL is
\[
 \varphi_n(x,y)=(L_y\sqrt\pi)^{-1/2}
      e^{\ii\kappa ny}e^{-(x-\kappa n)^2/2},\qquad
 \psi_j=\sum_{m\in\mathbb Z}\varphi_{j+mM},
 \quad 0\le j<M.
\]
The sums and their derivatives converge locally uniformly.
They satisfy \eqref{eq:bc} and the LLL equation.
Integration in \(y\) kills unequal Fourier modes, and the
remaining intervals in \(x\) tile the real line; hence the
\(\psi_j\) are orthonormal.  Conversely, Fourier expansion in
\(y\) and the LLL equation force a Gaussian in \(x\) for each
Fourier coefficient, while the \(x\)-boundary condition makes
the coefficients periodic modulo \(M\).  This proves completeness.

The projector kernel is
\(\Pi(r,r')=\sum_j\psi_j(r)\overline{\psi_j(r')}\).
Expanding the two sums and applying Poisson summation to the
remaining Fourier label expresses it as a sum indexed by
\((\ell_1,\ell_2)\in\mathbb Z^2\), whose summands have modulus
\[
 \frac1{2\pi}
 \exp\!\left[-\frac{(x-x'+\ell_1L_x)^2+(y-y'+\ell_2L_y)^2}{4}\right].
\]
For example, before Poisson summation the expression is
\[
 \frac1{L_y\sqrt\pi}\sum_{\ell_1\in\mathbb Z}
 e^{-\ii\ell_1L_xy'}e^{-(x-x'+\ell_1L_x)^2/4}
 \sum_{n\in\mathbb Z}e^{\ii\kappa n(y-y')}
 e^{-(\kappa n-(x+x'-\ell_1L_x)/2)^2}.
\]
The Poisson formula
\[
 \sum_n e^{-(\kappa n-a)^2}e^{\ii\kappa nb}
 =\frac{\sqrt\pi}{\kappa}\sum_{\ell_2}
       e^{-(b+\ell_2L_y)^2/4}e^{\ii a(b+\ell_2L_y)}
\]
gives the claimed modulus term by term.
Covariant derivatives of the kernel introduce polynomials in
the image displacement.  The connection terms cancel the
coordinate-dependent part of the magnetic phase.  In particular,
for a fixed number \(s\) of covariant derivatives,
\begin{equation}\label{eq:kernelbound}
 \abs{\nabla_{\mathrm{cov}}^{\,s}\Pi(r,r')}
 \le C_s\sum_{\lambda\in L_x\mathbb Z\oplus L_y\mathbb Z}
 (1+\abs{r-r'+\lambda})^s
                  e^{-\abs{r-r'+\lambda}^2/4}.
\end{equation}
The derivatives here have the appropriate dual connection
in the conjugate variable.

For \(L\ge4\), the image lattice has uniform separation.
If \(a_\lambda=\abs{r-R+\lambda}\), then
\(\sum_\lambda e^{-a_\lambda^2/16}\le C\).
Cauchy--Schwarz, with part of the Gaussian as a weight, therefore
bounds the square of the right side of \eqref{eq:kernelbound}
by \(C_s\sum_\lambda e^{-a_\lambda^2/8}\), after absorbing
polynomials into the Gaussian.
On the set \(d(r,R)>r_0\) all \(a_\lambda>r_0\).
Integrating over a fundamental domain and summing over images
is bounded by the plane integral
\[
 C_s\int_{\abs z>r_0}e^{-\abs z^2/8}\dd z
 \le C_s'e^{-c r_0^2}.
\]
This proves the \(L^2\) tail bounds for \(e_R=\Pi(\cdot,R)\)
and for \(g_R\), which is obtained by differentiating the
evaluation functional.  Their uniform norm bounds follow
by setting \(r_0=0\).  Finally, the orthonormal expansion gives
\[
 \int |e_R\rangle\langle e_R|\dd R
 =\sum_{j,k}|\psi_j\rangle\langle\psi_k|
       \int\overline{\psi_j(R)}\psi_k(R)\dd R=\Pi.
\]

\section{Locality without a localized LLL basis}\label{app:locality}

\subsection{Spatial cells and a bounded interaction}

Partition the physical torus into rectangles of side lengths
between \(1\) and \(2\).  Their adjacency graph has uniformly
bounded degree and quadratic volume growth.  A cell carries
the infinite-dimensional Fock space of its \(L^2\) sections;
no basis in that space is chosen.

Replace every test vector in \(k_R\) by its restriction to the
ball of radius \(r\) about \(R\).  The resulting \(k_R^{(r)}\)
is bounded, neutral, and supported in that ball.  A telescopic
expansion of a product of annihilators and creators, together
with \(\norm{a(f)}=\norm f\), gives
\begin{equation}\label{eq:kapprox}
 \norm{k_R}\le C,\qquad
 \norm{k_R-k_R^{(r)}}\le Ce^{-cr^2}.
\end{equation}
The degrees and coefficients are fixed, so the constants are
uniform.  The truncation respects adjoints.  We take its neutral
part by averaging over the two global gauge rotations if necessary;
this preserves support and does not increase the error, since
\(k_R\) is neutral.  Equivalently one can first express \(k_R\)
as a sum of monomials neutral in each spin and then truncate.

For each cell \(z\), put \(A_z=\int_z k_R\dd R\), and similarly
define \(A_z^{(r)}\).  Increasing the radius by a fixed harmless
amount, define
\[
 \Phi_{z,0}=A_z^{(1)},\qquad
 \Phi_{z,n}=A_z^{(n+1)}-A_z^{(n)}\quad(n\ge1).
\]
Then \(H_b=\sum_{z,n}\Phi_{z,n}\), each term is supported on
cells within distance \(C(n+1)\) of \(z\), and
\(\norm{\Phi_{z,n}}\le Ce^{-cn^2}\).
On a finite torus the truncations eventually equal \(A_z\).
This is a decomposition by physical regions, not an expansion
in one-particle orbitals.

If terms with the same support are collected into \(\Phi(Z)\),
the preceding bounds imply, for any fixed \(a>0\),
\begin{equation}\label{eq:Fbound}
 \mathcal I_a:=\sup_{x,y}
 \frac{\displaystyle\sum_{Z\ni x,y}\norm{\Phi(Z)}}
      {F_a(d(x,y))}<\infty,\qquad
 F_a(r)=\frac{e^{-ar}}{(1+r)^4},
\end{equation}
uniformly in the torus.  Indeed a contributing radius is at
least a constant times \(d(x,y)\), the possible centers grow
only quadratically, and a Gaussian absorbs every such factor.
Quadratic volume growth also gives
\[
 \sum_z F_a(d(x,z))F_a(d(z,y))
       \le C_F F_a(d(x,y)).
\]
To see this, split the sum according to which of
\(d(x,z),d(z,y)\) is at least \(d(x,y)/2\), use the triangle
inequality for the exponentials, and sum the remaining
\((1+r)^{-4}\) factor.

\subsection{Propagation estimates}

\begin{lemma}[Lieb--Robinson estimate]\label{lem:LR}
For bounded even ambient operators \(A,B\) supported on cell
sets \(X,Y\), with \(X\cap Y=\varnothing\),
\begin{equation}\label{eq:LR}
 \norm{[\tau_t(A),B]}
 \le C\norm A\norm B
       \sum_{x\in X,y\in Y}F_a(d(x,y))(e^{v\abs t}-1),
 \qquad \tau_t(A)=e^{\ii tH_b}Ae^{-\ii tH_b}.
\end{equation}
The constants are independent of the dimension of the cell
Hilbert spaces and of \(L\).
\end{lemma}
\begin{proof}
Set
\[
 C_B(X,t)=\sup_{A\in\mathcal A_X^{\rm even},\,\norm A\le1}
                         \norm{[\tau_t(A),B]}.
\]
It suffices to consider \(t\ge0\): replacing \(H_b\) by
\(-H_b\) treats negative times with the same interaction norms.
For a local observable supported on a set \(X\), pass to the
interaction picture with respect to the terms contained in
\(X\).  Differentiating its commutator with \(B\), and using
Duhamel's formula, bounds the propagation out of \(X\) by
twice the sum of the norms of interaction terms meeting \(X\),
integrated against the corresponding commutator bound for
their supports.  Enlarging the sum to all terms meeting \(X\),
the resulting inequality is
\[
 C_B(X,t)\le C_B(X,0)
 +2\sum_{Z:Z\cap X\ne\varnothing}\norm{\Phi(Z)}
                \int_0^t C_B(Z,s)\dd s.
\]
Here \(C_B(Z,0)\le2\norm B\), and it vanishes if \(Z\cap Y\)
is empty.  Iterating this integral inequality gives
chains \(Z_1,\ldots,Z_n\) with
\[
 Z_1\cap X\ne\varnothing,\quad
 Z_j\cap Z_{j+1}\ne\varnothing,\quad
 Z_n\cap Y\ne\varnothing.
\]
Their ordered time integrals are \(\abs t^n/n!\).
Choose a site at each intersection and use \eqref{eq:Fbound}.
The sum of chain weights is at most
\[
 \mathcal I_a^n C_F^{\,n-1}
             \sum_{x\in X,y\in Y}F_a(d(x,y)).
\]
Summing the series with the commutator factor \(2^n\) proves
\eqref{eq:LR}, with \(v=2\mathcal I_aC_F\) and an enlarged constant.
All steps use norms, unitary conjugation, and commuting
disjoint even algebras; no trace on a local space is used.
The finite-volume interaction sum is norm convergent, so
finite-sum approximations justify all differentiations.
\end{proof}

We also need a local approximation to the dynamics.  Let
\(H_{X^{(r)}}\) contain only interaction terms whose supports
lie in the \(r\)-neighborhood of \(X\), and let
\(\tau_t^{X^{(r)}}\) denote its dynamics.  Duhamel's formula gives
\[
 \tau_t(A)-\tau_t^{X^{(r)}}(A)
 =\ii\int_0^t \tau_s\!
 \left([H_b-H_{X^{(r)}},
              \tau_{t-s}^{X^{(r)}}(A)]\right)\dd s.
\]
Let \(\Lambda_L\) denote the cell graph.  An omitted term of
diameter less than \(r/2\) is at distance at least \(r/2\)
from \(X\).  Apply \eqref{eq:LR} and use
\(\sum_Z|Z|\norm{\Phi(Z)}\le C|\Lambda_L|\), which follows
from \eqref{eq:Fbound} with equal sites.  For terms of larger
diameter the shell decomposition gives
\(\sum_{\diam Z\ge r/2}\norm{\Phi(Z)}
 \le C|\Lambda_L|(1+r)^c e^{-cr^2}\).
Write \(I_{\mathrm{short}}(u)\) and \(I_{\mathrm{long}}(u)\)
for the corresponding sums of commutator norms in the Duhamel
integrand, with evolved observable \(\tau_u^{X^{(r)}}(A)\).
The preceding estimates give explicitly
\begin{align*}
 I_{\mathrm{short}}(u)
 &\le C\norm A\,|X|\,|\Lambda_L|e^{-ar/2}
                        (e^{v\abs u}-1),\\
 I_{\mathrm{long}}(u)
 &\le C\norm A\,|\Lambda_L|(1+r)^c e^{-cr^2}.
\end{align*}
The second line uses only \(\norm{[B,C]}\le2\norm B\norm C\).
Integrating in time, absorbing the factor \(\abs t\) into a
larger exponential, and reducing the spatial decay rate gives,
for some \(a',v'>0\),
\begin{equation}\label{eq:dynamicsapprox}
 \norm{\tau_t(A)-\tau_t^{X^{(r)}}(A)}
 \le C\norm A\,|X|\,|\Lambda_L|(1+r)^c e^{v'\abs t-a'r}.
\end{equation}
The same estimate holds for restricted interactions and is
uniform in \(L\).  Its volume prefactor is harmless because
\(|\Lambda_L|=O(L^2)\), and we use it at \(r\) proportional
to \(L\).  The truncated dynamics preserves neutrality.

\subsection{Boundary currents and unbounded charges}

On the finite-particle core,
\[
 [Q_\Gamma,a_\up(f)]=-a_\up(1_\Gamma f),\qquad
 [Q_\Gamma,a_\up^\dagger(f)]=a_\up^\dagger(1_\Gamma f).
\]
There is no spin-down contribution.  These formulas show
directly that \(j_R=\ii[k_R,Q_\Gamma]\) extends to a bounded
ambient operator.  If \(R\) is inside \(\Gamma\), replace
\(1_\Gamma f\) by \(f\) in each factor.  The resulting
commutator vanishes because each monomial is neutral.
Only tails outside \(\Gamma\) remain.  For \(R\) outside
\(\Gamma\), all \(1_\Gamma f\) are tails.  Thus
\begin{equation}\label{eq:currentdecay}
 \norm{j_R}\le Ce^{-c\,d(R,\partial\Gamma)^2}.
\end{equation}
Replacing the test vectors by their restrictions to a ball
about \(R\) also gives local approximants to \(j_R\) with
Gaussian error.  This estimate follows from the displayed
commutator formulas, not from multiplying a norm error by
the unbounded ambient norm of \(Q_\Gamma\).

Partition the centers \(R\) according to the nearer boundary
component and integrate \(j_R\) over the two pieces.  This
defines \(J_\pm\), with
\[
 J_-+J_+=\ii[H_b,Q],\qquad \norm{J_\pm}\le CL_y.
\]
For a local approximant within transverse radius \(r\), omit
centers farther than \(r/2\) from the relevant boundary and
cut off their test vectors at radius \(r/2\).
Equation~\eqref{eq:currentdecay} and the Gaussian field tails
give an error at most \(CL_y(1+r)^c e^{-cr^2}\).
The identity with the commutator holds first on the core
and then on each fixed-number sector, where \(Q\) is bounded.

\section{Filters, dressing, and clustering in a sector}\label{app:filters}

\subsection{The inverse-energy filter}

Let \(\Delta=\min\{\gamma,\mu-C_{\mathrm{rem}}-w\}\) be the
lower bound on the gap above the band from
Lemma~\ref{lem:generalcomparison}.  The choice
\(\mu>C_{\mathrm{rem}}+w+\gamma\) in
Proposition~\ref{prop:analytic} gives \(\Delta=\gamma>2w\).
For a positive parent with lowest energy zero,
\(C_{\mathrm{rem}}=0\).  On the selected number sector,
\(H_b\) has the same band projection, width, and gap, with
band energies in \([E_1-\mu N,E_d-\mu N]\).
Choose a real smooth even function \(\chi\) equal to one
on a neighborhood of \([-w,w]\) and zero for
\(\abs\omega\ge\Delta\).
With Fourier convention
\(\widehat W(\omega)=\int W(t)e^{\ii t\omega}\dd t\), choose
\[
 \widehat W(\omega)=\frac{1-\chi(\omega)}{\ii\omega}.
\]
Its inverse Fourier transform is real, bounded, and integrable,
with
\begin{equation}\label{eq:Wmoments}
 \int(1+\abs t)^m\abs{W(t)}\dd t<\infty
 \quad\hbox{for every }m.
\end{equation}
For an explicit construction, let
\[
 \check\chi(t)=\frac1{2\pi}\int\chi(\omega)e^{-\ii t\omega}\dd\omega.
\]
This is a real even Schwartz function with
\(\int\check\chi(t)\dd t=1\).
For \(t\ne0\), set
\begin{equation}\label{eq:Wexplicit}
 W(t)=-\operatorname{sgn}(t)\int_{\abs t}^{\infty}\check\chi(s)\dd s,
 \qquad W(0)=0.
\end{equation}
It is real, odd, and bounded, with one-sided limits
\(W(0^\pm)=\mp1/2\).  Its tail integral decays faster than
every power, proving \eqref{eq:Wmoments}.  Distributionally,
\(W'=\check\chi-\delta_0\), so
\(-\ii\omega\widehat W(\omega)=\chi(\omega)-1\).
This proves the asserted Fourier formula, with its continuous
extension at \(\omega=0\); oddness gives \(\widehat W(0)=0\).

Define bounded ambient operators
\[
 K_\pm^\infty=\int_{\mathbb R}W(t)\tau_t(J_\pm)\dd t,\qquad
 K^\infty=K_-^\infty+K_+^\infty.
\]
The integrals converge in norm.  On the fixed-number sector,
matrix elements between the low-energy band and its orthogonal
complement have energy difference of magnitude at least \(\Delta\).
For such a difference \(\omega\),
\(\widehat W(\omega)\ii\omega=1\).
The two off-diagonal blocks relative to \(P\) therefore give
\begin{equation}\label{eq:dressingexact}
 [K^\infty,P]=[Q,P].
\end{equation}
This can equivalently be checked using spectral measures;
no finite-dimensionality of the enlarged sector is required.
The energy differences inside the band lie in \([-w,w]\),
where \(\widehat W=0\).  Their individual values do not enter
\eqref{eq:dressingexact}, which only concerns the two
off-diagonal blocks relative to \(P\).

To localize \(K_\pm^\infty\), approximate \(J_\pm\) in a strip
of radius \(r/4\), use the restricted dynamics in a strip of
radius \(r/2\), and integrate only over \(\abs t\le c_0r\).
Choose \(c_0\) so that \(v'c_0<a'/8\) in
\eqref{eq:dynamicsapprox}, allowing for the distance \(r/4\)
between the two strip edges.  The errors for these times are
exponentially small in \(r\), with polynomial boundary
prefactors.  The omitted times are controlled by
\eqref{eq:Wmoments}.  For every \(m\) this gives a
self-adjoint neutral approximant \(K_\pm^{(r)}\) satisfying
\begin{equation}\label{eq:Kapprox}
 \supp K_\pm^{(r)}\subset\{d(\cdot,\partial_\pm)<r\},\qquad
 \norm{K_\pm^\infty-K_\pm^{(r)}}
       \le C_m L^c(1+r)^{-m}.
\end{equation}
These are ambient norm bounds.  Take \(r=L_x/64\) to obtain
the \(K_\pm\) in Proposition~\ref{prop:analytic}.
Equation~\eqref{eq:dressingexact} then gives the required
approximate commutation with \(P\).

\subsection{Clustering across the low-energy projection}

We give a filter proof that uses a gap only in the selected
sector.  Choose \(\vartheta\in C^\infty(\mathbb R)\) with
\[
 \vartheta(\omega)=1\quad(\omega\le-\Delta),\qquad
 \vartheta(\omega)=0\quad(\omega\ge-\Delta/2).
\]
For a bounded neutral local \(A\), write
\(A_{\vartheta}=\vartheta(\ad_{H_b})A\).  This is unambiguous in finite volume:
\(H_b\) acts on the finite-dimensional LLL Fock factor and as
the identity on its complement, so it has finitely many distinct
eigenvalues.  If \(H_b=\sum_\lambda\lambda E_\lambda\), set
\[
 A_{\vartheta}=\sum_{\lambda,\lambda'}
 \vartheta(\lambda-\lambda')E_\lambda A E_{\lambda'}.
\]
For matrix elements within \(\Ran P\), the energy difference
belongs to \([-w,w]\), where \(\vartheta=0\) because
\(w<\Delta/2\).  For a matrix element with output in
\(\Ran P\) and input in its complement, the energy difference
is at most \(-\Delta\), where \(\vartheta=1\).  For the
reverse matrix element it is at least \(\Delta\), where
\(\vartheta=0\).  Neutrality restricts these statements to
the selected number sector.  Thus, even with distinct energies
inside the band, the exact identities are
\[
 PA_{\vartheta}=PA(1-P),\qquad A_{\vartheta}P=0.
\]
Consequently, for another bounded neutral local \(B\),
\begin{equation}\label{eq:clusteridentity}
 PA(1-P)BP=P[A_{\vartheta},B]P.
\end{equation}
The identities are sector identities; \(\vartheta\) need not remove
transitions in other charge sectors.

For clarity, the distributional Fourier representation of
this filter introduces no norm singularity into the commutator.
Since \(\vartheta'\) is smooth and compactly supported, the inverse
transform of \(\vartheta\) is a multiple of \(\delta(t)\) plus
\(b(t)/t\) in principal value, with \(b\) a Schwartz function.
For disjoint supports the delta term vanishes in
\([A_{\vartheta},B]\), and the integral at \(t=0\) converges because
\eqref{eq:LR} contains \(e^{v\abs t}-1=O(\abs t)\).

Let \(s=d(X,Y)\).  Split this integral at \(\abs t=\eta s\),
where \(v\eta<a/2\).  The short-time contribution is bounded
by an exponentially small quantity times
\(\norm A\norm B|X||Y|\), using \eqref{eq:LR}.
For the long-time contribution use
\(\norm{[\tau_t(A),B]}\le2\norm A\norm B\) and the rapid
decay of \(b(t)/t\).  Thus, for every \(m\),
\begin{equation}\label{eq:clusterbound}
 \norm{PA(1-P)BP}
 \le C_m\norm A\norm B\,|X||Y|(1+s)^{-m}.
\end{equation}
The estimate also holds for local von Neumann algebra
operators, such as regional-charge exponentials: the
commutator proof of Lemma~\ref{lem:LR} applies directly to
bounded operators on the cell Fock spaces.
For the macroscopic separated strips \(S_\pm\), their cell
counts are polynomial in \(L\), so \eqref{eq:clusterbound}
is precisely \eqref{eq:groundcluster}.

\subsection{Bookkeeping for sector-bounded charges}

The only operators in the index proof that lack a bounded
ambient norm are regional charges, the corresponding
translation transports, and their sums with bounded dressings.
Their restrictions to \(\HH_{\bm n}\) have polynomial norms.
All errors involving those operators are estimated after
restriction, for example
\[
 \norm{[Q_X,E]|_{\HH_{\bm n}}}
 \le2n_\up\norm{E|_{\HH_{\bm n}}}
\]
for neutral \(E\).
Their exponentials, on the other hand, are ambient local
unitaries of norm one; these are the operators to which the
clustering and separated-factor arguments are applied.
Bounded perturbations of a local charge are self-adjoint on
its domain and have local unitary exponentials.
Thus no step applies a Lieb--Robinson estimate with an
unbounded ambient observable, and no step regards a global
particle-number compression as a local operation.

\begin{thebibliography}{9}
\bibitem{BBDF19}
S.~Bachmann, A.~Bols, W.~De Roeck, and M.~Fraas,
\emph{A many-body index for quantum charge transport},
Commun.\ Math.\ Phys.\ \textbf{375}, 1249--1272 (2020),
\href{https://arxiv.org/abs/1810.07351}{arXiv:1810.07351v3}.

\bibitem{BBDF21}
S.~Bachmann, A.~Bols, W.~De Roeck, and M.~Fraas,
\emph{Rational indices for quantum ground state sectors},
J.\ Math.\ Phys.\ \textbf{62}, 011901 (2021),
\href{https://arxiv.org/abs/2001.06458v3}{arXiv:2001.06458v3}.
Version 3 (31 March 2022) corrects Lemma 4.1.

\bibitem{Haldane85}
F.~D.~M.~Haldane,
\emph{Many-Particle Translational Symmetries of Two-Dimensional
Electrons at Rational Landau-Level Filling},
Phys.\ Rev.\ Lett.\ \textbf{55}, 2095--2098 (1985),
\href{https://doi.org/10.1103/PhysRevLett.55.2095}
{doi:10.1103/PhysRevLett.55.2095}.
\end{thebibliography}
\end{document}
